% Options for packages loaded elsewhere
\PassOptionsToPackage{unicode}{hyperref}
\PassOptionsToPackage{hyphens}{url}
\documentclass[
  a4paper,
]{article}
\usepackage{xcolor}
\usepackage[margin=25mm]{geometry}
\usepackage{amsmath,amssymb}
\setcounter{secnumdepth}{-\maxdimen} % remove section numbering
\usepackage{iftex}
\ifPDFTeX
  \usepackage[T1]{fontenc}
  \usepackage[utf8]{inputenc}
  \usepackage{textcomp} % provide euro and other symbols
\else % if luatex or xetex
  \usepackage{unicode-math} % this also loads fontspec
  \defaultfontfeatures{Scale=MatchLowercase}
  \defaultfontfeatures[\rmfamily]{Ligatures=TeX,Scale=1}
\fi
\usepackage{lmodern}
\ifPDFTeX\else
  % xetex/luatex font selection
\fi
% Use upquote if available, for straight quotes in verbatim environments
\IfFileExists{upquote.sty}{\usepackage{upquote}}{}
\IfFileExists{microtype.sty}{% use microtype if available
  \usepackage[]{microtype}
  \UseMicrotypeSet[protrusion]{basicmath} % disable protrusion for tt fonts
}{}
\makeatletter
\@ifundefined{KOMAClassName}{% if non-KOMA class
  \IfFileExists{parskip.sty}{%
    \usepackage{parskip}
  }{% else
    \setlength{\parindent}{0pt}
    \setlength{\parskip}{6pt plus 2pt minus 1pt}}
}{% if KOMA class
  \KOMAoptions{parskip=half}}
\makeatother
\usepackage{longtable,booktabs,array}
\newcounter{none} % for unnumbered tables
\usepackage{calc} % for calculating minipage widths
% Correct order of tables after \paragraph or \subparagraph
\usepackage{etoolbox}
\makeatletter
\patchcmd\longtable{\par}{\if@noskipsec\mbox{}\fi\par}{}{}
\makeatother
% Allow footnotes in longtable head/foot
\IfFileExists{footnotehyper.sty}{\usepackage{footnotehyper}}{\usepackage{footnote}}
\makesavenoteenv{longtable}
\setlength{\emergencystretch}{3em} % prevent overfull lines
\providecommand{\tightlist}{%
  \setlength{\itemsep}{0pt}\setlength{\parskip}{0pt}}
\usepackage{newunicodechar}
\newunicodechar{↔}{\ensuremath{\leftrightarrow}}
\newunicodechar{→}{\ensuremath{\rightarrow}}
\newunicodechar{←}{\ensuremath{\leftarrow}}
\newunicodechar{≥}{\ensuremath{\ge}}
\newunicodechar{≤}{\ensuremath{\le}}
\newunicodechar{×}{\ensuremath{\times}}
\newunicodechar{≪}{\ensuremath{\ll}}
\newunicodechar{≫}{\ensuremath{\gg}}
\newunicodechar{∈}{\ensuremath{\in}}
\newunicodechar{∞}{\ensuremath{\infty}}
\newunicodechar{±}{\ensuremath{\pm}}
\newunicodechar{−}{\ensuremath{-}}
\newunicodechar{≈}{\ensuremath{\approx}}
\newunicodechar{≠}{\ensuremath{\ne}}
\newunicodechar{∝}{\ensuremath{\propto}}
\newunicodechar{√}{\ensuremath{\sqrt{\phantom{x}}}}
\newunicodechar{⊕}{\ensuremath{\oplus}}
\newunicodechar{α}{\ensuremath{\alpha}}
\newunicodechar{β}{\ensuremath{\beta}}
\newunicodechar{θ}{\ensuremath{\theta}}
\newunicodechar{ρ}{\ensuremath{\rho}}
\newunicodechar{π}{\ensuremath{\pi}}
\newunicodechar{φ}{\ensuremath{\varphi}}
\newunicodechar{λ}{\ensuremath{\lambda}}
\newunicodechar{σ}{\ensuremath{\sigma}}
\newunicodechar{Δ}{\ensuremath{\Delta}}
\newunicodechar{Σ}{\ensuremath{\Sigma}}
\newunicodechar{ν}{\ensuremath{\nu}}
\newunicodechar{κ}{\ensuremath{\kappa}}
\newunicodechar{μ}{\ensuremath{\mu}}
\newunicodechar{τ}{\ensuremath{\tau}}
\newunicodechar{ω}{\ensuremath{\omega}}
\newunicodechar{ε}{\ensuremath{\varepsilon}}
\newunicodechar{Ψ}{\ensuremath{\Psi}}
\newunicodechar{⁻}{\ensuremath{{}^{-}}}
\newunicodechar{¹}{\ensuremath{{}^{1}}}
\newunicodechar{²}{\ensuremath{{}^{2}}}
\newunicodechar{³}{\ensuremath{{}^{3}}}
\newunicodechar{⋆}{\ensuremath{\star}}
\newunicodechar{∗}{\ensuremath{*}}
\newunicodechar{★}{\ensuremath{\bigstar}}
\newunicodechar{∎}{\ensuremath{\blacksquare}}
\newunicodechar{ℝ}{\ensuremath{\mathbb{R}}}
\newunicodechar{ℤ}{\ensuremath{\mathbb{Z}}}
\newunicodechar{₀}{\ensuremath{{}_{0}}}
\newunicodechar{₁}{\ensuremath{{}_{1}}}
\newunicodechar{₂}{\ensuremath{{}_{2}}}
\newunicodechar{₃}{\ensuremath{{}_{3}}}
\newunicodechar{₄}{\ensuremath{{}_{4}}}
\newunicodechar{₅}{\ensuremath{{}_{5}}}
\newunicodechar{₆}{\ensuremath{{}_{6}}}
\newunicodechar{₇}{\ensuremath{{}_{7}}}
\newunicodechar{₈}{\ensuremath{{}_{8}}}
\newunicodechar{₉}{\ensuremath{{}_{9}}}
\newunicodechar{✓}{\ensuremath{\checkmark}}
\newunicodechar{—}{---}
\newunicodechar{①}{(1)}
\newunicodechar{②}{(2)}
\newunicodechar{③}{(3)}
\usepackage{bookmark}
\IfFileExists{xurl.sty}{\usepackage{xurl}}{} % add URL line breaks if available
\urlstyle{same}
\hypersetup{
  hidelinks,
  pdfcreator={LaTeX via pandoc}}

\author{}
\date{}

\begin{document}

\section{Structure of Counting Readouts in an Anonymous Two-Channel
Closed Wave
System}\label{structure-of-counting-readouts-in-an-anonymous-two-channel-closed-wave-system}

\subsection{Exact Rationality of Equal-Weight Points, the Meta Counting
Rule, the Convention Audit, Complete Verification of the Niven
Intersection Points, and Exact Lock States Constructed by Counting
Alone}\label{exact-rationality-of-equal-weight-points-the-meta-counting-rule-the-convention-audit-complete-verification-of-the-niven-intersection-points-and-exact-lock-states-constructed-by-counting-alone}

\textbf{Author:} Noriaki Kihara\\
\textbf{ORCID:} 0009-0004-6753-4020\\
\textbf{Date:} August 3, 2026\\
\textbf{Version DOI:} \texttt{10.5281/zenodo.21763998}\\
\textbf{Concept DOI:} \texttt{10.5281/zenodo.21763997}\\
\textbf{Position:} ``Wave Information Readout'' series, two-channel
exchange scattering system, supplementary paper v1 (ninth-paper series;
\textbf{Part II of a trilogy}. Part I = Two-Grammar Decomposition
{[}4{]}, Part III = Lock Dynamics {[}5{]})

\begin{center}\rule{0.5\linewidth}{0.5pt}\end{center}

\subsection{Abstract}\label{abstract}

\textbf{(Decomposition theorem)} The readout of an anonymous two-channel
closed wave system decomposes exactly into ``counting + amplitude
deformation.'' At the point where the harmonics are given equal weights
(the equal-weight point), the readout degenerates to a counting of the
existence set of harmonics and becomes an \textbf{exact rational number
with a small denominator} (examples: \(61/128\) for odd 1..63, \(29/64\)
for odd 1..31, \(15/31\) when 3 is removed, exactly \(0\) for the even
control), and its value can be predicted solely from the measured bin
table of each harmonic and the accounting of the mask (\textbf{meta
counting rule}: error \(\le8\times10^{-16}\) across all 20/20 conditions
of 7 held-out supports × 5 convention variants, with the predictions
fixed as in-code constants before measurement). Amplitude distributions,
on the other hand, contaminate the readout continuously (spread
\(0.407\) for an identical existence set). \textbf{The only carrier of
structural information is the existence set of harmonics.}

\textbf{(Convention audit)} This structure survives across 5 readout
convention variants (changes of the mask definition). What does not
survive --- the concrete form of the counting rule, the concrete values
of the equal-weight points, and the role assignment of odd/even --- is
bookkeeping. In particular, under the parity-inverted mask, the support
that carries the lock moves to \textbf{30 even harmonics}: ``odd
harmonics = fermionic'' is a convention label, and the structure
(counting → lock) is independent of the label.

\textbf{(Complete verification of the Niven intersection points)} The
intersection of the rational numbers \(R\) reachable by counting with
rational-angle locking \(\theta/\pi\in\mathbb{Q}\) is, by Niven's
theorem, strictly limited to the 5 points
\(R\in\{0,\,1/4,\,1/2,\,3/4,\,1\}\). At every reachable intersection
point we verified the direct link ``counting construction (equal
weights, \textbf{zero amplitude search}) → exact lock,'' with all
predictions fixed before measurement, and scored \textbf{6/6 hits}: the
mixed-parity support \(\{2,5\}\) and its kin land exactly on \(R=1/4\)
with period 12 and \((m,n)=(1,3)\); the \(\{1,3\}\)-free support gives
\(R=1/2\), period 8, \((1,4)\); odd 5..63 under the variant mask gives
\(R=3/4\), period 6, \((1,6)\); and \(R=0\) is invariant (all errors
\(\le3.9\times10^{-16}\), recurrence residuals \(\sim10^{-15}\)).
\textbf{The realized orders \(n=3,4,6\) coincide exactly with the
complete menu of nontrivial orders of the two-dimensional
crystallographic restriction} --- Niven's theorem and the
crystallographic restriction are two faces of the same rational-cosine
constraint.

\textbf{(Where the integers live, and the limit)} The integer pair
\((m,n)\) can be read out from the order and winding number of the
evolution (measured: exact reconstruction in all 3 state families). By
contrast, a static state decomposition carries no winding number:
winding-number spectroscopy on a cyclic register showed that the state
spectrum coincides completely with mapping artifacts and that the
occupancy of the target winding number is at random level (a refutation
experiment). Hence the generation rule for integer pairs
\textbf{outside} the 5 Niven points (including the physical
\(\,(23,124)\)) can only be sought in dynamics, not in static structure
--- that verdict is delivered by Part III {[}5{]}.

\textbf{Keywords:} counting readout, rational degeneration, meta
counting rule, convention audit, Niven's theorem, crystallographic
restriction, exact recurrence, reproducible computation

\begin{center}\rule{0.5\linewidth}{0.5pt}\end{center}

\subsection{0. Conclusion}\label{conclusion}

\[
\boxed{
\begin{aligned}
&\text{Readout = counting + amplitude deformation. Only the existence set carries structure (exact rationals at equal-weight points, }\le8\times10^{-16}\text{ on 20/20).}\\
&\text{Role labels (odd = fermionic, etc.) are convention; the structure is label-independent (under parity inversion the even side locks).}\\
&\text{The direct counting}\to\text{lock link scores full hits at the 5 Niven intersection points, with realized orders }\{3,4,6\}\text{ = the crystallographic menu.}\\
&\text{Integer pairs can be read only from dynamics — static decomposition carries no winding number (refuted).}
\end{aligned}
}
\]

\subsection{Position of This Paper (Three-Part
Structure)}\label{position-of-this-paper-three-part-structure}

Part I {[}4{]} treats the binary identity decomposition of the flow and
the two grammars; this paper (Part II) treats the structural layer of
the readout --- what is counting and what is appearance; Part III
{[}5{]} treats the dynamics (the reality of locking and the limits of
selection). The three parts share a single reproduction package.

\subsection{1. Research Question}\label{research-question}

How much of a readout is the structure of the system, and where does the
product of the way of reading begin? From where can the integers be
read? {[}4{]} established the compositional grammar of readouts, but the
status of the readout \textbf{value} itself --- the boundary between the
structural layer (independent of convention) and the bookkeeping layer
(a product of convention) --- remained unclassified. This paper
completes that classification experimentally and verifies the full range
of lock states that counting can build (the Niven intersection points).

\subsection{2. System Definition and Design
Boundaries}\label{system-definition-and-design-boundaries}

The system, state construction, standard readout, and rotation are
identical to Section 2 of {[}4{]}. In all experiments the scattering
core and the readout formula are unmodified, target values are used only
for state construction, predictions are fixed as in-code constants
before measurement, and each runner begins with a reproduction of known
results (an anchor).

\subsection{3. Experiment I: Two-Memory Battery (Discovery of the
Decomposition)}\label{experiment-i-two-memory-battery-discovery-of-the-decomposition}

\textbf{Condition A (identical existence set, different amplitude
distributions, same pair norm)}: fixing the support to odd 1..63 and
varying the amplitude distribution over 5 types (equal weights, linear
decay, \(1/k\), exponential decay, random), the readout \(R\) moves
across \(0.069\)--\(0.477\) (\textbf{spread \(0.407\)}). A
power-ratio-type readout is strongly contaminated by amplitude (the
visible layer).

\textbf{Condition B (same norm, different existence sets, equal
weights)}: at the equal-weight point the readout degenerates to an exact
rational number.

{\def\LTcaptype{none} % do not increment counter
\begin{longtable}[]{@{}ll@{}}
\toprule\noalign{}
Support & \(R\) (exact value) \\
\midrule\noalign{}
\endhead
\bottomrule\noalign{}
\endlastfoot
Odd 1..63 (32 harmonics) & \(61/128\) \\
Odd 1..31 (16 harmonics) & \(29/64\) \\
Odd step 4 (16 harmonics) & \(15/32\) \\
Odd with 3 removed (31 harmonics) & \(15/31\) \\
Even control & \(0\) (exact) \\
\end{longtable}
}

The denominator tracks the support size, and the exact zero of the even
control is a structural readout of parity. Control test: the
equal-weight single-harmonic sum agrees exactly with the template
construction (\(1.7\times10^{-16}\)). Invariance under evolution (100
collisions) has drift \(\le3.3\times10^{-16}\).

\textbf{Consequence: readout = counting readout + amplitude deformation.
The equal-weight point is the structural representative of each support
class.}

\subsection{4. Experiment II: The Meta Counting Rule and Held-Out
Prediction
Verification}\label{experiment-ii-the-meta-counting-rule-and-held-out-prediction-verification}

The counting rule decoded from the 5 points of Condition B (its concrete
form under the canonical mask),

\[
R_{\mathrm{eq}}(S)=\frac{2|S|-c(S)}{4|S|},\qquad
c:\{k{=}1\to2,\ k{=}3\to1,\ \text{even }k\to2,\ \text{odd }k{\ge}5\to0\}
\]

was verified as a \textbf{prediction fixed before measurement} on 7
held-out supports (\(\{1,3\}\), \(\{5\}\), odd 5..63, \(\{1,5,7\}\),
\(\{3,7\}\), odd 1..15, even-mixed \(\{1,2,5\}\)), and every one was a
hit (error \(\le5\times10^{-16}\), including the generalized \(c\) rule
for the even-mixed case).

However, the concrete form \(c\) is bookkeeping of the canonical mask
(Section 5). The convention-independent core is the \textbf{meta
counting rule} --- ``the \(R\) of an equal-weight point can be predicted
solely from the measured bin table of each harmonic alone and the
accounting of the mask'' --- and this survives under all conventions.

\subsection{5. Experiment III: Convention
Audit}\label{experiment-iii-convention-audit}

We swept the fermion mask over 5 variants (canonical even≥4 / even≥6 /
even≥2 / any≥4 / parity-inverted odd≥3) and measured all 20 combinations
with 4 state classes (a control test confirmed that the variant readout
agrees exactly with the original implementation in the canonical
setting, difference \(0.0\)).

\textbf{Physics-grade (survives all conventions)}: the meta counting
rule (20/20, \(\le8\times10^{-16}\)), the rationality of the
equal-weight points (\(61/128\), \(59/128\), \(95/128\), \(61/124\),
\(3/4\), \(1/4\), etc.), the amplitude dependence (excluding degenerate
classes where the readout is identically zero), and the counting → lock
bridgehead. \textbf{In particular, under the parity-inverted mask, 30
even harmonics lock} (period 8, \((1,4)\)) --- the convention relativity
of role assignment.

\textbf{Bookkeeping-grade (products of convention)}: the concrete form
of the \(c\) rule, the concrete values of the equal-weight points, the
role assignment of odd/even, and the existence condition of the
bridgehead (under even≥2, A is no longer purely bosonic and the counting
lock disappears).

All claims are stated in meta form (physics-grade). All audit
measurements are collected in a figure (\texttt{audit\_figure\_v1}).

\subsection{6. Experiment IV: Complete Verification of the Niven
Intersection
Points}\label{experiment-iv-complete-verification-of-the-niven-intersection-points}

\subsubsection{6.1 Intersection Theorem}\label{intersection-theorem}

The equal-weight points of counting give rational \(R\). On the other
hand, exact locking requires \(\theta/\pi\in\mathbb{Q}\) (a rational
angle). Compatibility of \(R=\cos^2\theta\) is, by Niven's theorem (if
\(\theta/\pi\in\mathbb{Q}\) and \(\cos\theta\in\mathbb{Q}\), then
\(\cos\theta\in\{0,\pm\tfrac12,\pm1\}\) {[}8{]}), limited via the
rationality of \(\cos2\theta\) to the 5 points

\[
R\in\{0,\ 1/4,\ 1/2,\ 3/4,\ 1\}
\]

The corresponding orders are \(n\in\{3,4,6\}\) (the nontrivial points)
--- \textbf{the complete menu of nontrivial orders of the
two-dimensional crystallographic restriction} (two faces of the same
rational-cosine constraint).

\subsubsection{6.2 Full Prediction Verification (6/6
Hits)}\label{full-prediction-verification-66-hits}

All predictions were fixed as in-code constants before measurement, and
every reachable intersection point was verified:

{\def\LTcaptype{none} % do not increment counter
\begin{longtable}[]{@{}
  >{\raggedright\arraybackslash}p{(\linewidth - 10\tabcolsep) * \real{0.1667}}
  >{\raggedright\arraybackslash}p{(\linewidth - 10\tabcolsep) * \real{0.1667}}
  >{\raggedright\arraybackslash}p{(\linewidth - 10\tabcolsep) * \real{0.1667}}
  >{\raggedright\arraybackslash}p{(\linewidth - 10\tabcolsep) * \real{0.1667}}
  >{\raggedright\arraybackslash}p{(\linewidth - 10\tabcolsep) * \real{0.1667}}
  >{\raggedright\arraybackslash}p{(\linewidth - 10\tabcolsep) * \real{0.1667}}@{}}
\toprule\noalign{}
\begin{minipage}[b]{\linewidth}\raggedright
Case
\end{minipage} & \begin{minipage}[b]{\linewidth}\raggedright
Support
\end{minipage} & \begin{minipage}[b]{\linewidth}\raggedright
Readout
\end{minipage} & \begin{minipage}[b]{\linewidth}\raggedright
\(R\) (prediction / error)
\end{minipage} & \begin{minipage}[b]{\linewidth}\raggedright
Period (prediction / measured)
\end{minipage} & \begin{minipage}[b]{\linewidth}\raggedright
\((m,n)\) (prediction / measured)
\end{minipage} \\
\midrule\noalign{}
\endhead
\bottomrule\noalign{}
\endlastfoot
Anchor & Odd 5..63 & Canonical & \(1/2\) / \(3.9\times10^{-16}\) & 8 /
\textbf{8} & \((1,4)\) / \textbf{\((1,4)\)} \\
P1a--c & Mixed parity \(\{2,5\}\) and 2 kin (3 types) & Canonical &
\(1/4\) / \(\le3.3\times10^{-16}\) & 12 / \textbf{12} & \((1,3)\) /
\textbf{\((1,3)\)} \\
P2 & Odd 5..63 & even≥2 variant & \(3/4\) / \(3.3\times10^{-16}\) & 6 /
\textbf{6} & \((1,6)\) / \textbf{\((1,6)\)} \\
P3 & \(\{1\}\) & Canonical & \(0\) / \(0.0\) & Invariant ✓ & --- \\
\end{longtable}
}

All of these are pure counting constructions with equal weights and
\textbf{zero amplitude search} (\(R=1\) is unreachable by counting
because the A fundamental wave is bosonic under all masks. Period-4
recurrence in the amplitude-tuned family has been confirmed in the
existing R-sweep experiments --- we record this honestly). The
mixed-parity supports (the P1 series) are a state class opened for the
first time by this verification.

\subsection{\texorpdfstring{7. Experiment V: A Reader for the Integer
Pair
\((m,n)\)}{7. Experiment V: A Reader for the Integer Pair (m,n)}}\label{experiment-v-a-reader-for-the-integer-pair-mn}

For an exact lock state, we measure the order (the collision count \(P\)
of the first exact recurrence) and the winding number (the accumulated
rotation angle over one period \(/2\pi\)), and obtain \((m,n)\) by
rational reconstruction of \(x=1/2-\bar\theta/\pi\). In all lock state
families the result agreed exactly with the prediction: \((1,3)\),
\((1,4)\), \((1,6)\) (winding-number error \(\le6\times10^{-16}\)).
Integer pairs can be read out from the evolution.

\subsection{8. Control and Refutation: Static Decomposition Carries No
Winding
Number}\label{control-and-refutation-static-decomposition-carries-no-winding-number}

On a cyclic register of non-common harmonics 31 × 4 phases = 124 slots,
we measured the winding-number spectrum of the antisymmetric component
of the state under 4 canonical orderings (winding-number spectroscopy).
The spectrum of the equal-weight state coincided completely with the
``mapping-only control'' under all orderings, every observed winding
structure (\(k=31,1,93\)) was an artifact of the slot mapping, and the
occupancy of the target winding number \(k=23\) was at random level
under all conditions. \textbf{A static initial-state decomposition
carries no winding number} (recorded as a refutation experiment).

Therefore, the generation rule for integer pairs outside the 5 Niven
points --- including the physical \((23,124)\) {[}2{]} --- can only be
sought in dynamics. The verdict on the dynamics side (the reality of
locking and the limits of selection) is delivered by Part III {[}5{]}.

\subsection{9. Claims}\label{claims}

\textbf{Claim 1 (Decomposition theorem).} The readout decomposes exactly
into counting (the structural layer: a function of the existence set,
exactly rational at equal-weight points, predictable by the meta
counting rule) and amplitude deformation (the visible layer: spread
0.407). The only carrier of structural information is the existence set
of harmonics.

\textbf{Claim 2 (Convention relativity).} The structural layer survives
across 5 readout convention variants, and the role labels (odd =
fermionic, etc.) are convention. Structure alone remains without labels
--- the readout-convention version of the anonymity principle of
{[}4{]}.

\textbf{Claim 3 (Niven intersection points).} The direct link counting →
exact lock holds at the 5 Niven intersection points, and only there. The
realized orders \(\{3,4,6\}\) are the complete menu of the
crystallographic restriction. All predictions scored 6/6 hits under
pre-measurement fixing.

\textbf{Claim 4 (Where the integers live).} The integer pair \((m,n)\)
can be read out from the evolution (order and winding number) and cannot
be read out from static decomposition (refuted). The generation rule for
general \((m,n)\) is restricted to the dynamical type.

\subsection{10. Open Problems}\label{open-problems}

\begin{enumerate}
\def\labelenumi{\arabic{enumi}.}
\tightlist
\item
  \textbf{A closed-form general theory of the counting rule}:
  generalization of the meta counting rule (bin table × mask) to
  arbitrary state constructions and arbitrary carrier structures
\item
  \textbf{Outside the 5 Niven points}: dynamical generation of
  \((23,124)\)-type integer pairs --- already transformed, by the
  elimination argument of Part III {[}5{]}, into the question of ``the
  commutativity of the observation clock''
\item
  \textbf{The \(R=1\) intersection point}: counting reachability depends
  on the construction of the A anchor. A general theory including the
  anchor degree of freedom
\end{enumerate}

\subsection{11. Reproducibility}\label{reproducibility}

We follow the same discipline as {[}4{]}. All runners, all measurement
CSVs, and figures are committed:

{\def\LTcaptype{none} % do not increment counter
\begin{longtable}[]{@{}
  >{\raggedright\arraybackslash}p{(\linewidth - 4\tabcolsep) * \real{0.3333}}
  >{\raggedright\arraybackslash}p{(\linewidth - 4\tabcolsep) * \real{0.3333}}
  >{\raggedright\arraybackslash}p{(\linewidth - 4\tabcolsep) * \real{0.3333}}@{}}
\toprule\noalign{}
\begin{minipage}[b]{\linewidth}\raggedright
Experiment
\end{minipage} & \begin{minipage}[b]{\linewidth}\raggedright
Folder
\end{minipage} & \begin{minipage}[b]{\linewidth}\raggedright
Commit
\end{minipage} \\
\midrule\noalign{}
\endhead
\bottomrule\noalign{}
\endlastfoot
Two-memory battery & \texttt{two\_memory\_readout\_battery\_pre\_v1} &
5444fbdb \\
Counting rule verification + lock link + \((m,n)\) readout &
\texttt{counting\_rule\_lock\_bridge\_pre\_v1} & 8dff131e \\
Convention audit & \texttt{convention\_audit\_pre\_v1} & 5a8f4b0a \\
Audit figure & same folder, \texttt{run\_audit\_figure\_v1} &
47e2da02 \\
Niven intersection-point lock & \texttt{niven\_points\_lock\_pre\_v1} &
8f99d5c6 \\
Winding-number spectroscopy (refutation) &
\texttt{winding\_spectroscopy\_pre\_v1} & 641048ad \\
\end{longtable}
}

\begin{center}\rule{0.5\linewidth}{0.5pt}\end{center}

\section{References}\label{references}

\subsection{Self-citations}\label{self-citations}

\begin{enumerate}
\def\labelenumi{\arabic{enumi}.}
\tightlist
\item
  Noriaki Kihara, ``Anonymous Equal-Amplitude Composite-Wave Model:
  Basic Axiom System v6,'' Version DOI:
  \texttt{10.5281/zenodo.21465984}, Concept DOI:
  \texttt{10.5281/zenodo.21315735}, 2026.
\item
  Noriaki Kihara, ``Discovery of Finite-Order Resonance in Iterated
  Exchange Scattering,'' Version DOI: \texttt{10.5281/zenodo.21421367},
  Concept DOI: \texttt{10.5281/zenodo.21421366}, 2026.
\item
  Noriaki Kihara, ``Future Phase-Position Acceleration Map and the
  Inverse-Square Law via Harmonic Closure in an AB Two-Body Closed Phase
  System v4,'' Version DOI: \texttt{10.5281/zenodo.21468270}, 2026.
  (Precedent for the three-step procedure)
\item
  Noriaki Kihara, ``Two-Grammar Decomposition of Interaction in an
  Anonymous Two-Channel Closed Wave System'' (Part I of the trilogy),
  Version DOI: \texttt{10.5281/zenodo.21763996}, Concept DOI:
  \texttt{10.5281/zenodo.21763995}, 2026.
\item
  Noriaki Kihara, ``Lock Dynamics in an Anonymous Two-Channel Closed
  Wave System'' (Part III of the trilogy), Version DOI:
  \texttt{10.5281/zenodo.21764000}, Concept DOI:
  \texttt{10.5281/zenodo.21763999}, 2026.
\item
  Noriaki Kihara, ``Exchange-Interference Scattering Matrix \ldots{}
  Summary of Preliminary Locality-Exchange Experiments v1,'' Version
  DOI: \texttt{10.5281/zenodo.21333768}, 2026. (Code provenance of the
  engine)
\item
  Noriaki Kihara, ``Numerical Experiments on Exchange-Scattering
  Coefficient R Concentration and Fine-Structure-Constant Correspondence
  Candidates v1,'' Version DOI: \texttt{10.5281/zenodo.21396761}, 2026.
  (Code provenance of System A)
\end{enumerate}

\subsection{External References}\label{external-references}

\begin{enumerate}
\def\labelenumi{\arabic{enumi}.}
\setcounter{enumi}{7}
\tightlist
\item
  I. Niven, \emph{Irrational Numbers}, Carus Mathematical Monographs
  No.~11, MAA (1956). (The theorem on the rationality of cosines of
  rational angles. The common root of the intersection set and the
  crystallographic restriction.)
\end{enumerate}

\begin{center}\rule{0.5\linewidth}{0.5pt}\end{center}

\textbf{Addendum (Record of a Correction)} In the course of this
research, an error was once recorded that took the intersection of the
counting rational family and the rational-angle lock family to be
``\(R=1/2\) only'' (both the occurrence and the correction are appended
and preserved in the experiment folder README). The correction to the 5
points by Niven's theorem is Section 6 of this paper, and the
verification of the remaining 2 points arising from the correction (P1
and P2) was completed with all predictions hitting. We record, together
with its history, how the error led to the discovery of a new state
class (the mixed-parity supports).

\end{document}
